\documentclass[12pt]{article}
\usepackage[T1]{fontenc}
\usepackage[utf8]{inputenc}
\usepackage{lmodern}
\usepackage{textcomp}
\usepackage{enumitem}
\usepackage{geometry}
\geometry{margin=1in}
\begin{document}
\begin{center}
Answer Key - Political Science - Graduate Level Assessment 25 \
\small (Answers are ordered exactly as in the question paper.)
\end{center}

\section*{Multiple Choice}
\begin{enumerate}[label=\arabic*.]
\item \textbf{Answer:} C
\\ \textit{Options:} A) containment.; B) neoconservatism.; C) isolationism.; D) protectionism.
\\ \textit{Note:} Isolationism is the correct answer because, for much of American history, particularly prior to World War II, the United States largely avoided entanglement in foreign conflicts and prioritized domestic affairs over international involvement.
\item \textbf{Answer:} B
\\ \textit{Options:} A) All of the world's leading economies were declining due to low growth and inflation; B) The United States could no longer remain a superpower and was in decline; C) The soft power of the United States would allow it to avoid decline; D) The rise of Japan had been exaggerated
\\ \textit{Note:} Paul Kennedy argued that the United States was experiencing a decline as a superpower due to the overextension of its military commitments and economic challenges, which ultimately undermined its global influence and power.
\item \textbf{Answer:} D
\\ \textit{Options:} A) Limiting security analysis to military engagement is too restrictive.; B) Security is a contestable concept.; C) Focusing security only on violent conflicts is too restrictive.; D) All of these options.
\\ \textit{Note:} The correct answer is "All of these options" because the Human Development Report has significantly broadened the scope of human security to encompass various interconnected dimensions such as economic stability, environmental sustainability, and social inclusion, making all listed ideas integral to the mainstream understanding of human security.
\item \textbf{Answer:} D
\\ \textit{Options:} A) the president of the United States.; B) the secretary of defense.; C) the chairman of the Joint Chiefs of Staff.; D) Congress.
\\ \textit{Note:} The power to declare war rests with Congress as outlined in Article I, Section 8 of the U.S. Constitution, which grants this authority explicitly to the legislative branch, ensuring a system of checks and balances between the branches of government.
\item \textbf{Answer:} A
\\ \textit{Options:} A) There is no specific referent point.; B) The state and state apparatus.; C) The environment.; D) The people and their collectives.
\\ \textit{Note:} In Critical Security Studies, the absence of a specific referent object emphasizes the perspective that security is a fluid and socially constructed concept, focusing on broader issues such as human security and structural inequalities rather than solely state-centric paradigms.
\end{enumerate}

\section*{True / False}
\begin{enumerate}[label=\arabic*.]
\setcounter{enumi}{5}
\item \textbf{Answer:} False
\\ \textit{Options:} A) True; B) False
\\ \textit{Note:} The statement is false because contemporary cyber-security frameworks prioritize the protection of individuals, organizations, and critical infrastructure over the mere digitalization of sensitive information, recognizing that these entities are the primary referent objects at risk.
\item \textbf{Answer:} True
\\ \textit{Options:} A) True; B) False
\\ \textit{Note:} The statement is true because global health organizations, including the World Health Organization and UNAIDS, report that around 38 million people were living with HIV/AIDS as of the latest estimates, confirming the figure falls within the stated range.
\item \textbf{Answer:} False
\\ \textit{Options:} A) True; B) False
\\ \textit{Note:} The 'New Populism' is characterized by a nationalist and often anti-globalization stance, emphasizing sovereignty and local interests over international cooperation and alliances.
\item \textbf{Answer:} False
\\ \textit{Options:} A) True; B) False
\\ \textit{Note:} Idealism emphasizes the importance of various actors, including international organizations and non-state entities, in international relations, and it advocates for moral principles and cooperation among these diverse participants rather than solely focusing on states.
\item \textbf{Answer:} False
\\ \textit{Options:} A) True; B) False
\\ \textit{Note:} Most states maintain secrecy around their nuclear capabilities to deter adversaries and protect national security, which undermines the idea of promoting transparency and international trust.
\end{enumerate}

\section*{Long Form Response}
\begin{enumerate}[label=\arabic*.]
\setcounter{enumi}{10}
\item \textbf{Answer:} Inefficient balancing or buckpassing by states can significantly contribute to a more competitive international system by creating power vacuums and fostering uncertainty among states. For instance, when a state fails to adequately counter a rising power, it may encourage aggressive behavior from that power, as seen in the pre-World War I context where major European powers neglected to confront Germany's expansionist policies. This lack of decisive action can lead to a security dilemma, where states feel compelled to enhance their military capabilities, thereby escalating tensions and competition. Additionally, buckpassing-where states rely on others to manage threats-can result in weaker alliances and a fragmented international response to crises, exemplified by the failure to effectively address the Syrian civil war. Ultimately, these dynamics not only heighten competition among states but also complicate cooperative efforts to address global challenges, thereby reshaping the landscape of international politics.
\\ \textit{Note:} correct because it illustrates how states' failure to effectively balance against rising powers can lead to increased competition and insecurity in the international system, as demonstrated by historical examples like the pre-World War I dynamics involving Germany.
\item \textbf{Answer:} The differing relationships between U.S. and European security studies can be attributed to historical, geopolitical, and cultural factors. Historically, the U.S. has positioned itself as a global hegemon with a focus on military intervention and a liberal international order, while European security studies often emphasize diplomacy, multilateralism, and regional stability, reflecting their experiences with war and integration. Geopolitically, the U.S. prioritizes counter-terrorism and great power competition, particularly with nations like Russia and China, whereas Europe tends to focus on internal security and transnational threats such as migration and climate change. Culturally, differing attitudes toward state sovereignty and intervention influence their security policies. These variations have significant implications for international relations, as they shape alliances, influence policy-making, and impact cooperative efforts in addressing global security challenges.
\\ \textit{Note:} The answer correctly identifies that the differences in U.S. and European security studies stem from their distinct historical experiences, geopolitical priorities, and cultural attitudes, which shape their approaches to international relations and security policy.
\end{enumerate}

\vfill
\noindent\textit{End of Answer Key}
\end{document}